% Exercises for Chapter 10 (stratified sampling and conditional Monte Carlo).
% Statements only: solutions are kept privately.
% Unless stated otherwise: the test case S = K = 100, T = 1, r = 1%, y = 2%, sigma = 20%,
% S_T = S exp(nu T + sigma sqrt(T) Z) with nu = r - y - sigma^2/2 and Z ~ N(0,1).

\begin{exercise}[Stratifying a digital]\label{ex:st-digital}
\easy Let $D=e^{-rT}\1_{\{S_T>K\}}$ in the Black--Scholes model, $z^*=(\ln(K/S)-\nu T)/(\sigma\sqrt T)$,
$u^*=N(z^*)$ and $p=1-u^*=\Q(S_T>K)$. Stratify $U=N(Z)$ into $H$ equiprobable strata with
proportional allocation, $n$ a multiple of $H$. You may use \noteref{thm:st-variance} and
\noteref{prop:st-optimal}.
\begin{enumerate}[label=(\alph*)]
  \item Show that at most one stratum has $\sigma_h>0$, and that
    $n\Var\hat\mu^{\mathrm{st}}_n=e^{-2rT}\vartheta(1-\vartheta)/H$, where $\vartheta\in[0,1]$
    is the fraction of that stratum lying above $u^*$. Deduce
    $n\Var\hat\mu^{\mathrm{st}}_n\le e^{-2rT}/(4H)$ and compare with plain sampling.
  \item With $n\ge H$ and at least one draw in every stratum, what allocation minimises the
    variance, and what is the minimal variance?
  \item When is the stratified estimator exact? What does this suggest about where to put the
    edges of the strata for a payoff with a jump?
  \item For $K=110$, $H=10$ and $n=1000$, give the standard-error ratios to plain sampling of
    the estimators in (a) and (b).
\end{enumerate}
\end{exercise}

\begin{exercise}[Allocation with costs]\label{ex:st-costs}
\easy Suppose that a draw in stratum $h$ costs $c_h>0$ (for instance, stratum $h$ is sampled
by acceptance--rejection with acceptance probability $a_h$, so that $c_h$ is proportional to
$1/a_h$), and that the budget is $C=\sum_hc_hn_h$.
\begin{enumerate}[label=(\alph*)]
  \item Treating the $n_h>0$ as real numbers, show that $\sum_hp_h^2\sigma_h^2/n_h$ is
    minimised under the budget constraint by $n_h$ proportional to $p_h\sigma_h/\sqrt{c_h}$,
    and that the minimal variance is $\big(\sum_hp_h\sigma_h\sqrt{c_h}\big)^2/C$.
  \item A plain draw is obtained by first drawing the stratum with probabilities $p_h$ and then
    sampling in it, so it costs $\bar c=\sum_hp_hc_h$ on average. Show that the efficiency of
    the cost-optimal stratified estimator relative to plain sampling is at least $1$.
  \item Take two strata with $p=(0.9,0.1)$, $\sigma_h=(1,10)$, $c_h=(1,25)$ and $\mu_1=\mu_2$.
    Compare the products of variance and budget of plain sampling, proportional allocation,
    the Neyman allocation of \noteref{prop:st-optimal} (which ignores costs) and the
    cost-optimal allocation.
\end{enumerate}
\end{exercise}

\begin{exercise}[One point per stratum]\label{ex:st-one-point}
\medium Let $g:[0,1]\to\R$ and $\hat\mu_n=\frac1n\sum_{k=1}^ng\big((k-1+U_k)/n\big)$ with
$U_1,\dots,U_n$ i.i.d.\ uniform: $n$ equiprobable strata and one point in each. You may use
\noteref{thm:st-variance}.
\begin{enumerate}[label=(\alph*)]
  \item Let $g$ be continuously differentiable on $[0,1]$. Show that
    $n^3\Var\hat\mu_n\to\frac1{12}\int_0^1g'(u)^2\,\dd u$. \emph{Hint:} compare
    $g((k-1+U)/n)$ with its first-order expansion around $(k-1)/n$, and use the uniform
    continuity of $g'$.
  \item Deduce the rate of the root-mean-square error, and compare it with plain Monte Carlo and
    with \noteref{thm:st-variance}(v).
  \item For the discounted test-case call as a function of the uniform,
    $g(u)=e^{-rT}(S_T(N^{-1}(u))-K)^+$, show that $\int_0^1g'(u)^2\,\dd u=\infty$.
  \item This part is computational. For that call, compute $n^2\Var\hat\mu_n$ and
    $n^3\Var\hat\mu_n$ for $n=2^4,2^6,\dots,2^{16}$ (exactly, from the stratum variances, or
    by simulation with independent replications), and describe the rate you observe. Which
    strata dominate the variance?
\end{enumerate}
\end{exercise}

\begin{exercise}[Collapsed strata]\label{ex:st-collapsed}
\medium With one draw per stratum no within-stratum variance can be estimated
(\noteref{prop:st-error}). Let $H=2J$ equiprobable strata of a uniform $U$ receive one draw
$Y_h$ each, so that $\hat\mu^{\mathrm{st}}=\frac1H\sum_hY_h$, and pair the strata as
$(1,2),(3,4),\dots$.
\begin{enumerate}[label=(\alph*)]
  \item Show that $\hat v=\frac1{H^2}\sum_{j=1}^J(Y_{2j-1}-Y_{2j})^2$ satisfies
    $\E\hat v=\Var\hat\mu^{\mathrm{st}}+\frac1{H^2}\sum_j(\mu_{2j-1}-\mu_{2j})^2$: it is a
    conservative estimate of the variance.
  \item Let $Y=g(U)$ with $g$ continuously differentiable on $[0,1]$. Using the expansion of
    \exref{ex:st-one-point}(a), show that $\E\hat v/\Var\hat\mu^{\mathrm{st}}\to7$ as
    $H\to\infty$, unless $g$ is constant.
  \item The alternative of \noteref{prop:st-error}(iii) splits the $H$ draws into $R$
    independent designs of $H/R$ strata each. For the same $g$, compare the variance of the
    resulting estimator with that of the single design. Which way of obtaining an error
    estimate do you prefer, and why?
\end{enumerate}
\end{exercise}

\begin{exercise}[Post-stratification is a regression]\label{ex:st-post-regression}
\medium Let $(\xi_i,Y_i)$, $i\le n$, be i.i.d., with strata $A_1,\dots,A_H$ of known
probabilities $p_h>0$, and let $D^h_i=\1_{A_h}(\xi_i)$ and $N_h=\sum_iD^h_i$. You may use
\noteref{thm:cv-multiple}, \noteref{prop:cv-estimated-beta} and \noteref{prop:st-post}.
\begin{enumerate}[label=(\alph*)]
  \item Why can the $H$ indicators not all be used as controls together with an intercept?
    Use the $H-1$ controls $D^1,\dots,D^{H-1}$, with known means $p_1,\dots,p_{H-1}$.
  \item Show that, when every $N_h\ge1$, the regression estimator of
    \noteref{thm:cv-multiple}(iii) with these controls equals the post-stratified estimator
    $\sum_hp_h\bar Y_h$.
  \item Compute the $R^2$ of \noteref{thm:cv-multiple} for these controls, and check that
    $\sigma^2(1-R^2)$ is the asymptotic variance of \noteref{prop:st-post}(iv).
  \item An estimated control coefficient usually causes a bias of order $1/n$
    (\noteref{prop:cv-estimated-beta}). Why is the bias here exponentially small?
\end{enumerate}
\end{exercise}

\begin{exercise}[Neyman allocation for a deep out-of-the-money call]\label{ex:st-neyman-otm}
\medium This is a computational exercise. Take the course's deep out-of-the-money call
$S=100$, $K=150$, $T=1$, $r=5\%$, $y=0$, $\sigma=20\%$, with $n=10^5$ draws and $H$
equiprobable strata of $U=N(Z)$.
\begin{enumerate}[label=(\alph*)]
  \item For $H=100$, how many strata have $\mu_h=\sigma_h=0$? Explain why they can be left out
    of the simulation altogether, and write the resulting estimator.
  \item Compute the $\sigma_h$ of the other strata and the Neyman allocation of the $n$ draws
    among them. Give the variance factor $n\Var/\sigma^2$, and compare it with proportional
    allocation.
  \item Repeat for $H=10$ and $H=1000$, and explain the case $H=10$ in one line.
  \item Confirm (b) by simulation, with $100$ repetitions as in the course.
  \item Compare with importance sampling by a mean shift: draw $\tilde Z=Z+z^*$, with $z^*$ the
    value of $Z$ at which $S_T=K$, and average $f(\tilde Z)e^{-z^*\tilde Z+(z^*)^2/2}$, where
    $f(z)$ is the discounted payoff (you may use that this estimator is unbiased). Give its
    variance factor at equal $n$, and comment.
\end{enumerate}
\end{exercise}

\begin{exercise}[Latin hypercube sampling can lose]\label{ex:st-lhs-interaction}
\hard Let $V_1,\dots,V_n$ be a Latin hypercube sample (\noteref{def:st-lhs}) and
$f\in L^2([0,1]^d)$ with variance $\sigma^2$.
\begin{enumerate}[label=(\alph*)]
  \item Show that $\Var\hat\mu^{\mathrm{LHS}}_n=\frac{\sigma^2}n+\big(1-\frac1n\big)\Cov\big(f(V_1),f(V_2)\big)$.
  \item Let $d=2$ and $f(u)=(u_1-\frac12)(u_2-\frac12)$. Show that $\mu=0$, $\sigma^2=\frac1{144}$,
    and $\Cov(f(V_1),f(V_2))=c_n^2$ with $c_n=\E\big[(V_{11}-\frac12)(V_{21}-\frac12)\big]$.
    Show that $c_n=-\frac{n+1}{12n^2}$. \emph{Hint:} $V_{11}$ and $V_{21}$ lie in two distinct
    strata of the first coordinate, chosen uniformly at random.
  \item Deduce that $n\Var\hat\mu^{\mathrm{LHS}}_n/\sigma^2=1+(n-1)(n+1)^2/n^4>1$, so that Latin
    hypercube sampling is worse than plain sampling for this $f$, and check the bound
    $\sigma^2/(n-1)$ of \noteref{prop:st-lhs}(iv). How close to it is this example?
  \item Reconcile (c) with \noteref{prop:st-lhs}(iii).
  \item Check (c) by simulation for $n=4$ and $n=32$.
\end{enumerate}
\end{exercise}

\begin{exercise}[Conditioning a spread option]\label{ex:st-spread}
\medium This is a computational exercise. Under $\Q$ let
$S^i_T=S^i\exp\big((r-y_i-\frac12\sigma_i^2)T+\sigma_i\sqrt TZ_i\big)$, $i=1,2$, where $(Z_1,Z_2)$
is a standard normal pair with correlation $\rho\in(-1,1)$, and let
$Y=e^{-rT}(S^1_T-S^2_T-K)^+$ with $K\ge0$. You may use the Black--Scholes formulas and
\noteref{thm:st-conditional}.
\begin{enumerate}[label=(\alph*)]
  \item Writing $Z_1=\rho Z_2+\sqrt{1-\rho^2}\,Z^\perp$ with $Z^\perp$ standard normal and
    independent of $Z_2$, show that $\E[Y\mid Z_2]$ is a Black--Scholes call price, and
    identify its spot, strike and volatility.
  \item Take $S^1=S^2=100$, $K=5$, $T=1$, $r=1\%$, $y_1=y_2=2\%$, $\sigma_1=20\%$,
    $\sigma_2=30\%$, $\rho=0.5$. With $n=10^5$ draws, compare the plain and the conditional
    estimators: estimates, standard errors and the variance ratio.
  \item Condition instead on $Z_1$: show that $\E[Y\mid Z_1]$ is a put on $S^2_T$, and compare
    the two conditional estimators. Explain which one removes more variance.
  \item Stratify the conditioning variable of your better estimator into $100$ equiprobable
    strata. Which results of the chapter predict what you observe?
\end{enumerate}
\end{exercise}

\begin{exercise}[Hull and White's expansion]\label{ex:st-hull-white}
\hard Let $\rho=0$ in \noteref{prop:st-mixing}, with the Heston variance
$\dd V_t=\kappa(\theta-V_t)\,\dd t+\xi\sqrt{V_t}\,\dd W_t$ started at $V_0=\theta$, and let
$v=Q_T/T$ be the realised variance. Write $P(v)$ for the Black--Scholes put price as a function
of the variance $v=\sigma^2$, the other inputs being those of the test case. You may use that
$\Var V_s=\frac{\theta\xi^2}{2\kappa}(1-e^{-2\kappa s})$ and
$\Cov(V_s,V_t)=e^{-\kappa(t-s)}\Var V_s$ for $s\le t$.
\begin{enumerate}[label=(\alph*)]
  \item Show that $\E v=\theta$ and
    $\Var(Q_T)=\frac{\theta\xi^2}{\kappa^2}\Big(T-\frac{1-e^{-2\kappa T}}{2\kappa}-\frac{(1-e^{-\kappa T})^2}{\kappa}\Big)$.
  \item With $\mathcal V=Se^{-yT}N'(d_+)\sqrt T$ the Black--Scholes vega at volatility $\sqrt v$,
    show that $P''(v)=\mathcal V\,(d_+d_--1)/(4v^{3/2})$. Deduce the approximation
    $\E[P(v)]\approx P(\theta)+\frac12P''(\theta)\Var v$, and explain the sign of the correction
    at the money.
  \item Evaluate the approximation for $\kappa=2$, $\theta=0.04$, $\xi=0.3$, and compare it with
    the Fourier price $8.1534$ and with your own conditional Monte Carlo estimate (exact variance
    steps as in the chapter, or any scheme you trust).
  \item Repeat (c) with $\xi=0.6$ and $\xi=0.9$, using conditional Monte Carlo as the reference.
    Where does the expansion break down, and why?
\end{enumerate}
\end{exercise}
